\chapter{Módulo de Gestión} \label{chp:gestion}
Este primer módulo, denominado módulo de gestión, constituye el núcleo de la plataforma. Está dedicado íntegramente a la gestión de la aplicación, incluyendo la administración de usuarios, la autenticación y la arquitectura visual y de datos asociada al inicio del sistema.

\section{Gestión de Usuarios y Autenticación}

\subsection{Proceso de Login y Peticiones a Firebase}

Este módulo se encargará del login de los usuarios autorizados en mi empresa. Para agilizar y facilitar el proceso, el usuario deberá introducir tan solo su correo y contraseña para acceder a su homepage.

\subsubsection{Uso de Firebase Authentication} Usamos Firebase Authentication para gestionar el sistema de accesos de manera integral y segura. Cuando un usuario introduce su correo y contraseña, la aplicación cliente envía una petición de autenticación (\textit{signInWithEmailAndPassword}) a los servidores de Firebase. Si las credenciales son válidas, el servicio responde con un token cifrado de sesión y un identificador único (UID). Por otro lado, para registrar nuevos perfiles en la plataforma, el sistema lleva a cabo peticiones de creación de cuentas (\textit{createUserWithEmailAndPassword}), enlazando posteriormente esa identidad con el documento de permisos en la base de datos. También se gestionan a través de este servicio las peticiones para el restablecimiento de contraseñas olvidadas.

\subsubsection{Tipos de Usuarios} De entre los tres tipos de usuarios que se validan a través de este módulo, distinguimos los siguientes niveles de acceso y gestión: \begin{itemize} \item \textbf{Administrador:} El “admin” es el rol con mayores privilegios de la plataforma; podrá dar de alta tanto a empresas, como a usuarios, asignándoles el acceso correspondiente. \item \textbf{Recursos humanos:} El usuario “rrhh” tiene un perfil totalmente operativo; debería solo poder gestionar los usuarios de las empresas autorizadas a este, garantizando la privacidad de los datos. \item \textbf{Supervisor:} Considero que el supervisor debe ser un usuario que, dentro de la empresa, pueda dar de alta a más usuarios de rrhh. Esto permite delegar la administración del equipo interno a los responsables de las gestorías sin depender de un administrador global del sistema. \end{itemize}

\subsection{Delegación de Permisos} La jerarquía descrita mediante los distintos roles permite una delegación de permisos segura y estructurada. Una vez autenticado, el sistema reconoce inmediatamente el nivel de acceso del individuo y bloquea o permite las funcionalidades correspondientes para mantener un entorno de trabajo unificado pero completamente compartimentado.

\section{Estructura del Frontend}
El frontend de este módulo, implementado con tecnología \textit{Low-Code}, está estructurado principalmente en flujos de navegación que diferencian las áreas públicas de las privadas. La estructura del frontend se inicia en la \textbf{Login Page}, una vista sencilla que aloja el formulario de credenciales (correo y contraseña). No hay validaciones en tiempo real para el formato del correo o la fortaleza de la contraseña, ya que estas validaciones se delegan completamente a Firebase Authentication, que se encarga de verificar la existencia del usuario y la validez de las credenciales.

Una vez superado el login, el usuario es redirigido a la \textbf{Homepage}. La estructura interna de esta pantalla privada incluye un menú de navegación lateral (Drawer) fijo y un cuerpo central. La interfaz hace un fuerte uso de lógicas de \textbf{visibilidad condicional}; dependiendo del tipo de usuario (Admin, RRHH, Supervisor) que haya iniciado sesión, la estructura del frontend oculta o muestra automáticamente las distintas pantallas de gestión (como el formulario para dar de alta a una nueva empresa, que solo será visible para el administrador).

\section{Estructura del Backend}
El backend para este módulo funciona conectando Firebase Authentication con la base de datos documental (Firestore). Una vez se completa el registro o login de un usuario, el backend lleva a cabo la recuperación de un documento maestro en la colección \textit{users} que contiene la información completa del perfil, el rol asignado y la referencia a su empresa autorizada.
Para soportar las lógicas de segregación y seguridad (Multi-Tenant), en este apartado del backend cobra especial importancia el uso de \textbf{índices}. Dado que un usuario de "rrhh" solo debe acceder a los registros vinculados a la empresa que tiene autorizada, el sistema lleva a cabo consultas con filtros combinados (ej. buscar usuarios donde el campo \textit{empresa\_ref} sea X, y el campo \textit{rol} sea Y). Firestore exige la creación de \textbf{índices compuestos} específicos en el backend para poder estructurar previamente estos datos en memoria. La implementación de estos índices permite que las búsquedas y validaciones de permisos se ejecuten de manera casi instantánea y altamente eficiente, asegurando que nadie pueda realizar cruces de información no permitidos entre distintas empresas.